\begin{fragile}
    \frametitle{\insertsectionhead}
    \begin{block}{归并排序：分治策略的应用}
        将序列分为两半，分别递归排序，再将两个有序子序列合并。
    \end{block}
    \bicolumns[0.5]{
        \begin{block}{算法描述}
            \begin{algorithmx}[alg:merge_sort]{归并排序}
                \Procedure{MergeSort}{$s,l,r$}
                \State \Comment{排序左闭右开区间$[l,r)$}
                \If{$r-l\leq 1$}
                \State \textbf{return} $s[l:r]$
                \EndIf
                \State $m\gets l+\lfloor(r-l)/2\rfloor$
                \State $a\gets$ MergeSort($s,l,m$)
                \State $b\gets$ MergeSort($s,m,r$)
                \State \textbf{return} Merge($a,b$)
                \EndProcedure
            \end{algorithmx}
        \end{block}
    }{
        \begin{exampleblock}{实例}
            \begin{itemize}
                \item 原序列：$21,25,22,10,25,18$
                \item 分成：$(21,25,22)$、$(10,25,18)$
                \item 递归排序：$(21,22,25)$、$(10,18,25)$
                \item 合并结果：$10,18,21,22,25,25$
            \end{itemize}
        \end{exampleblock}
    }
\end{fragile}

\begin{frame}
    \frametitle{\insertsectionhead}
    \begin{block}{合并两个有序序列}
        \begin{itemize}
            \item 比较两个序列的当前元素，将较小者依次写入辅助数组
            \item 键相同时先取左侧元素，可保持相同键的原有相对顺序
            \item 一个序列耗尽后，复制另一个序列的剩余元素
        \end{itemize}
    \end{block}
    \begin{exampleblock}{实例：最后一次合并}
        $(21,22,25_a)$与$(10,18,25_b)$依次取出
        $10,18,21,22,25_a,25_b$；相等时先取左侧$25_a$。
    \end{exampleblock}
    \begin{alertblock}{性能}
        \begin{itemize}
            \item 时间：最好、平均和最坏均为$O(n\log n)$
            \item 常见数组实现的辅助空间为$O(n)$；上述合并规则保证稳定性
        \end{itemize}
    \end{alertblock}
\end{frame}

\begin{frame}
    \frametitle{\insertsectionhead}
    \begin{table}
        \centering
        \small
        \caption{常见数组排序算法比较}
        \begin{tabular}{lcccc}
            \toprule
            算法 & 最好时间 & 最坏时间 & 辅助空间 & 稳定 \\
            \midrule
            直接插入 & $O(n)$ & $O(n^2)$ & $O(1)$ & 是 \\
            直接选择 & $O(n^2)$ & $O(n^2)$ & $O(1)$ & 否 \\
            快速排序 & $O(n\log n)$ & $O(n^2)$ & $O(n)$ & 否 \\
            归并排序 & $O(n\log n)$ & $O(n\log n)$ & $O(n)$ & 是 \\
            \bottomrule
        \end{tabular}
    \end{table}
    \small 快速排序表中空间为所示递归实现的最坏情况；稳定性取决于具体实现。
\end{frame}

\begin{frame}
    \frametitle{\insertsectionhead：回到成绩记录}
    \begin{exampleblock}{选择稳定排序}
        对$21,25_a,22,10,25_b,18$采用示例中的稳定归并规则，
        得$10,18,21,22,25_a,25_b$；两位同分学生的原有次序得以保留。
    \end{exampleblock}
    \begin{block}{实际取舍}
        每次移动的是整条学生记录；归并排序提供$O(n\log n)$最坏时间，
        但常见数组实现需要$O(n)$辅助空间。
    \end{block}
\end{frame}

\begin{frame}{排序：思考与自测}
    \begin{itemize}
        \item 对$(2_a,1,2_b)$使用交换式直接选择排序，相同键的次序是否一定保持？
        \item 以成绩为主键、姓名为次键升序排列记录，如何利用稳定排序？
    \end{itemize}
    \pause
    \begin{block}<2-|handout:0>{参考答案}
        第一趟交换后为$(1,2_a,2_b)$；但$(2_a,2_b,1)$会变为
        $(1,2_b,2_a)$，故不稳定。多键排序先按姓名，再稳定地按成绩排序。
    \end{block}
\end{frame}
